\clearpage\begingroup\SourceLeaf{393b}\ArithmeticHead{62}{L}{STÄCKEL, GAUSS ALS GEOMETER.}
\renewcommand{\thefootnote}{\arabic{footnote})}\setcounter{footnote}{0}
\begin{center}
\textbf{16.}\par\smallskip
\textls[100]{Existenz der Wurzeln algebraischer Gleichungen.}
\end{center}
\noindent\hspace*{\parindent}Das Jahr 1797 brachte nicht nur die Entdeckungen über die lemniskatischen
Funktionen, damals ist auch der Beweis für die Existenz der Wurzeln
algebraischer Gleichungen entstanden, den \textsc{Gauss} in der Dissertation 1799
veröffentlicht hat (T. Nr. 80). Allerdings hat er es dort vermieden, imaginäre
Grössen zu benutzen. Schon im Titel hat er den zu beweisenden Satz in
der Form ausgesprochen, jede algebraische rationale ganze Funktion einer
Veränderlichen [mit reellen Koeffizienten] könne in reelle Faktoren ersten
oder zweiten Grades zerlegt werden, und im Art. 3 äussert er sich über die
imaginären Grössen in sehr vorsichtiger und zurückhaltender Weise. »Sollen
die imaginären Grössen überhaupt in der Analysis beibehalten werden, was
aus mehreren Gründen, die freilich hinreichend sichergestellt werden müssen,
richtiger scheint, als sie zu verwerfen, dann müssen sie notwendig für ebenso
möglich gelten wie die reellen . . . . Doch will ich mir die Rechtfertigung
der imaginären Grössen sowie eine eingehende Auseinandersetzung dieses ganzen
Gegenstandes für eine andere Gelegenheit vorbehalten« (W. III, S. 6).
\par
Dass \textsc{Gauss} damals schon im Besitze der geometrischen Versinnlichung
war, zeigt der Art. 16 (W. III, S. 22), denn die ganze Betrachtung läuft
darauf hinaus, dass die Funktion \(f(x+iy)\) in den reellen und den rein imaginären
Teil zerlegt wird und die Kurven in der \(xy\)-Ebene untersucht werden,
in denen je einer der beiden Teile verschwindet. Das sind die Spuren, die,
wie \textsc{Gauss} in der Selbstanzeige vom Jahre 1831 bemerkt hat, der aufmerksame
Leser in der Dissertation wiederfinden wird (W. II, S. 175). Hierzu ist freilich
zu bemerken, dass diese Andeutungen an und für sich nicht dazu ausreichen
würden, um den Schluss zu rechtfertigen, dass \textsc{Gauss} damals die geometrische
Versinnlichung der imaginären Grössen besessen habe, denn auch
\textsc{d’Alembert} hat in seinem Beweise für die Wurzelexistenz\footnote{J. \textsc{d’Alembert}, \textit{Recherches sur le calcul intégral}, 1. partie, Histoire de l’Acad. Année 1746, Berlin 1748, Mémoires, S. 182—191; vgl. P. \textsc{Stäckel}, \textit{Integration durch imaginäres Gebiet}, Bibliotheca math. (3) 1 (1900), S. 124.}, den \textsc{Gauss} im
\clearpage\SourceLeaf{394a}\ArithmeticHead{63}{R}{DIE KOMPLEXEN GRÖSSEN IN IHRER BEZIEHUNG ZUR GEOMETRIE.}\setcounter{footnote}{0}
\noindent Art. 5 wiedergibt und im Art. 6 beurteilt (W. III, S. 7—11), dasselbe Verfahren
benutzt, ohne dass ihm doch deshalb die geometrische Versinnlichung
der komplexen Grössen zuzuschreiben wäre.
\clearpage\endgroup

